\documentclass[12pt]{article}
\usepackage{amsmath}
\usepackage{amssymb}
\usepackage{amsthm}
\usepackage{geometry}
\geometry{margin=1in}
\usepackage{xcolor}
\usepackage{asymptote}
\usepackage{lastpage}
\usepackage{fancyhdr}
\pagestyle{fancy}
\fancyhf{}
\renewcommand{\headrulewidth}{0pt}
\cfoot{\thepage\ of \pageref{LastPage}}
\fancypagestyle{plain}{%
  \fancyhf{}%
  \renewcommand{\headrulewidth}{0pt}%
  \cfoot{\thepage\ of \pageref{LastPage}}%
}

\title{AMC 12 Prep 2026}
\author{}
\date{\today}

\setcounter{secnumdepth}{0}

\begin{document}
\maketitle

\section{2002 AMC12B Problem 10}

\textbf{Problem:}
How many different integers can be expressed as the sum of three distinct members of the set $\{1,4,7,10,13,16,19\}$ ?

$\text{(A)}\ 13 \qquad \text{(B)}\ 16 \qquad \text{(C)}\ 24 \qquad \text{(D)}\ 30 \qquad \text{(E)}\ 35$

\vspace{0.3cm}
\noindent\textbf{Answer:}~\underline{\hspace{5cm}}\hfill\textbf{Time:}~\underline{\hspace{1.2cm}}:\underline{\hspace{1.2cm}}

\vspace{0.5cm}

\section{2006 AMC12A Problem 10}

\textbf{Problem:}
For how many real values of $x$ is $\sqrt{120-\sqrt{x}}$ an integer?

$\textbf{(A) } 3\qquad \textbf{(B) } 6\qquad \textbf{(C) } 9\qquad \textbf{(D) } 10\qquad \textbf{(E) } 11$ .

\vspace{0.3cm}
\noindent\textbf{Answer:}~\underline{\hspace{5cm}}\hfill\textbf{Time:}~\underline{\hspace{1.2cm}}:\underline{\hspace{1.2cm}}

\vspace{0.5cm}

\end{document}